% 原书分章PDF第1—7页：章首及第15.1节，完整译文。
\chapter{聚集数据中的 LTMLE}
\label{ch:ltmle-clustering}
\chapterauthors{Mireille E. Schnitzer、Mark J. van der Laan、Erica E. M. Moodie 与 Robert W. Platt}
\chapteraffiliations{%
M. E. Schnitzer（通讯作者）：加拿大蒙特利尔大学药学院，\#2236, Pavillon Jean-Coutu 2940, chemin de la Polytechnique, Montreal, QC H3C 3J7, Canada；电子邮箱：\href{mailto:mireille.schnitzer@umontreal.ca}{mireille.schnitzer@umontreal.ca}。\\
M. J. van der Laan：美国加利福尼亚大学伯克利分校生物统计学部及统计学系，101 Haviland Hall, \#7358, Berkeley, CA 94720, USA；电子邮箱：\href{mailto:laan@berkeley.edu}{laan@berkeley.edu}。\\
E. E. M. Moodie 与 R. W. Platt：加拿大麦吉尔大学流行病学、生物统计学与职业卫生系，Purvis Hall, 1020 Pine Ave West, Montreal, QC H3A 1A2, Canada；电子邮箱：\href{mailto:erica.moodie@mcgill.ca}{erica.moodie@mcgill.ca}；\href{mailto:robert.platt@mcgill.ca}{robert.platt@mcgill.ca}。\\
本章 DOI：\url{https://doi.org/10.1007/978-3-319-65304-4_15}。}
\begingroup

母乳喂养被认为是婴儿早期喂养的最佳做法，也是大多数主要卫生组织所推荐的方式。然而，由于无法将母乳喂养直接作为随机干预进行分配，目前并没有直接的实验证据。母乳喂养促进干预试验（PROmotion of Breastfeeding Intervention Trial，PROBIT）是一项整群随机试验（cluster-randomized trial），旨在评价一项鼓励和支持母乳喂养的医院项目的效果，从而为母乳喂养预防婴儿感染和住院的保护作用提供间接证据。

\begin{sourcebox}
本章运用因果推断技术，估计不同母乳喂养持续时间这一纵向暴露，对出生后第一年内发生住院的时间区间数的影响。由于住院也可能影响母乳喂养能否继续，我们将其视为时变混杂因素（time-varying confounder）。我们将展示两种 g 计算方法，以及纵向靶向最小损失估计（longitudinal targeted minimum loss-based estimation，LTMLE）的一种实现方式；它们既考虑结局部分取决于时变混杂因素这一特点，也考虑研究设计本身造成的聚集性（clustering）。
\end{sourcebox}

\section{PROBIT 研究}

PROBIT 研究于1996年6月至1997年12月在白俄罗斯进行（Kramer 等，2001，2002）。研究将产科医院及其相应的综合门诊随机分配至接受或不接受哺乳管理培训，培训重点是如何鼓励延长纯母乳喂养。为尽可能提高效率，研究先按地区、城乡属性、每年分娩数以及出院时开始母乳喂养的比例，将医院匹配为17对，再在每对医院之间进行随机化。然而，由于两家医院退出，另有一家医院伪造记录，最终只有31个簇（cluster）完成研究，完整的配对结构也随之丧失。各医院簇仅招募打算对孩子进行母乳喂养的孕妇。具体而言，研究纳入的是健康、足月、单胎、出生体重不低于2500克的母乳喂养婴儿。基线数据包括母亲的人口学、教育和吸烟信息，既往妊娠情况，以及婴儿的信息：性别、出生体重、胎龄和 Apgar 评分（Finster 和 Wood，2005）。出生后第一年，在1、2、3、6、9和12个月进行随访，广泛收集婴儿喂养、生长、疾病和住院资料。31个簇共有17,044对母婴参与研究并有数据记录；其中8名研究对象缺少必需的基线数据，因此分析样本为17,036对母婴。表15.1列出分析中所调整的基线特征。

初步分析（Kramer 等，2001）发现，这项鼓励母乳喂养的试验对主要结局——胃肠道感染——有显著影响。随后，我们对母乳喂养持续时间影响第一年内胃肠道感染次数的效应进行了因果分析（Schnitzer 等，2014）。本章考察母乳喂养持续时间对婴儿住院的影响。我们尤其希望了解，延长母乳喂养是否会减少住院次数。由于问卷仅记录婴儿在两次访视之间是否住过院，因此12个月时的目标结局，是记录到住院的时间区间数。

\subsection{观测数据}

与 PROBIT 对应，我们考虑每对母婴的纵向数据，并将一对母婴定义为一名研究对象。数据形式为
\[
 O=(L_0,C_1,L_1,A_1,C_2,L_2,\ldots,L_{K-1},A_{K-1},C_K,Y),
\]
下标表示测量时点。$L_0$ 包含所有已测量的基线协变量，包括医院中心；$C_t$ 表示一次访视是否未能进行；$L_t$ 表示时变协变量；$A_t$ 是暴露水平的指示变量，如果婴儿在第 $t$ 个时间区间内一直接受母乳喂养，其取值为1；$Y$ 是第 $K$ 个时点测量的目标结局。图15.1表示时点 $t$ 所收集各项测量的先后顺序。对于任意时变变量，我们用 $\bar X_t$ 表示从开始直到并包括 $X_t$ 的历史。

\begin{table}[htbp]
\centering
\caption{PROBIT 数据集中17,044对母婴的基线特征}
\label{tab:ch15-baseline}
\small
\begin{tabular}{lrrr}
\toprule
特征 & \multicolumn{2}{c}{汇总} & 缺失数\\
\cmidrule(lr){2-3}
\textit{数值变量} & 中位数 & IQR\textsuperscript{a} & \\
\midrule
母亲年龄（岁） & 23 & $(21,27)$ & \\
既往子女数 & 0 & $(0,1)$ & \\
胎龄（个月） & 40 & $(39,40)$ & \\
婴儿体重（千克） & 3.4 & $(3.2,3.7)$ & \\
婴儿身长（厘米） & 52 & $(50,53)$ & \\
Apgar 评分\textsuperscript{b} & 9 & $(8,9)$ & 5\\
头围（厘米） & 35 & $(34,36)$ & 3\\
\addlinespace
\textit{二元变量} & 人数 & \% & \\
\midrule
妊娠期吸烟 & 389 & 2.28 & \\
过敏史 & 750 & 4.40 & \\
男婴 & 8827 & 52 & \\
剖宫产 & 1974 & 12 & \\
母亲教育程度 & & & 2\\
\quad 高中未毕业 & 663 & 4 & \\
\quad 高中毕业 & 5497 & 32 & \\
\quad 大学未毕业 & 8568 & 50 & \\
\quad 大学毕业 & 2316 & 14 & \\
地理区域 & & & \\
\quad 白俄罗斯东部，城市 & 5615 & 33 & \\
\quad 白俄罗斯东部，乡村 & 2706 & 16 & \\
\quad 白俄罗斯西部，城市 & 4380 & 26 & \\
\quad 白俄罗斯西部，乡村 & 4343 & 25 & \\
\bottomrule
\end{tabular}
\begin{minipage}{.96\textwidth}
\footnotesize
\textsuperscript{a}IQR：四分位距。\\
\textsuperscript{b}Apgar 评分用于评估新生儿健康状况，范围为1—10分；8分及以上表示状态良好，5—7分表示轻度抑制，4分及以下表示重度抑制（Finster 和 Wood，2005）。由于 PROBIT 的入组条件对基线体重和健康状况作了限制，实际观察到的评分范围为5—10分。本表及表题转载自 Schnitzer 等（2014）。
\end{minipage}
\end{table}

PROBIT 共有 $K=6$ 次随访。$Y$ 定义为第一年内婴儿至少发生过一次住院的时间区间数，因此是0至6之间的整数。我们关注的科学问题，是随时间测量的母乳喂养 $\{A_t;1\leq t\leq5\}$ 对 $Y$ 的影响。$L_t$ 表示婴儿是否在时间区间 $(t-1,t)$ 内住院。因此，若以 $L_6$ 表示第五与第六时点之间至少住院一次的指示变量，则 $Y=\sum_{t=1}^{6}L_t$。研究开始时，所有母亲都尝试母乳喂养，因此可以对所有研究对象定义 $A_0=1$。此外，我们只保留基线数据完整的子集，所以可以对所有研究对象定义 $C_0=0$，表示基线时未删失。

一旦某名研究对象错过一次访视，或者在访视时没有收集到必需的信息，我们就人为地将其在此后的所有访视中均视为删失。由于项目缺失并不常见，这种处理并未大幅减少可用数据。表15.2给出各时点删失、住院以及仍在母乳喂养的研究对象人数。注意，我们没有报告第6时点，即12个月时，仍在母乳喂养的人数，因为分析不使用这一信息。在从未删失的研究对象中，观察到发生住院的时间区间总数为4785。

\begin{figure}[htbp]
\centering
% 图15.1：按原图时间轴、括号、实心向下箭头及无箭头头的虚线连接重绘。
\begin{tikzpicture}[x=.55mm,y=-.55mm,>={Stealth[length=2.4mm,width=1.8mm]},font=\small]
 \draw (19.08,11.90) rectangle (131.04,29.57);
 \filldraw[fill=gray!22] (19.08,11.90) rectangle (36.75,29.57);
 \filldraw[fill=gray!22] (119.25,11.90) rectangle (131.04,29.57);
 \node at (27.92,20.73) {$t-1$};
 \node at (125.14,20.73) {$t$};
 \filldraw[fill=gray!22] (152.64,11.90) rectangle (228.67,29.57);
 \node at (190.66,20.73) {适用 $t=1,2,3,4,5$};
 \draw (36.75,32.52) .. controls (36.75,37.40) and (39.83,41.36) .. (43.63,41.36)
  -- (71.13,41.36) .. controls (74.92,41.36) and (78.00,45.31) .. (78.00,50.20)
  .. controls (78.00,45.31) and (81.08,41.36) .. (84.88,41.36)
  -- (112.37,41.36) .. controls (116.17,41.36) and (119.25,37.40) .. (119.25,32.52);
 \node[align=center,text width=44mm] at (78,72) {两次访视间隔内的住院：\\$L_t$};
 \draw[line width=1.0pt,->] (125.24,32.52) -- (125.24,91.45);
 \node[align=center,text width=45mm] at (130,113) {时点 $t$ 仍在母乳喂养\\（二元变量）：$A_t$};
 \draw[line width=1.0pt,dash pattern=on 1.0pt off 1.0pt]
  (39.00,25.29) .. controls (14.04,46.17) and (22.58,97.49) .. (58.07,139.93);
 \node[draw,line width=.6pt,densely dotted,anchor=north,minimum height=22.7mm,align=center,text width=45mm,inner sep=4pt]
  at (58,141) {时点 $t$ 的删失：$C_t$。\\错过时点 $t$ 的访视，\\意味着未测量 $L_t$ 和 $A_t$。};
\end{tikzpicture}

\caption{PROBIT 研究中变量的时间顺序。每次随访均记录母乳喂养状态（$A_t$）及过去一个时间区间内的住院情况（$L_t$）。在时点 $t$ 发生删失（$C_t=1$），表示没有观察到后续的母乳喂养及感染状态。}
\label{fig:ch15-time-order}
\end{figure}

\begin{table}[htbp]
\centering
\caption{各时点的删失人数、住院人数及仍在母乳喂养的母亲人数}
\label{tab:ch15-longitudinal}
\small
\setlength{\tabcolsep}{5pt}
\begin{tabular}{lrrrrrr}
\toprule
时点 & 1 & 2 & 3 & 4 & 5 & 6\\
月份 & 1 & 2 & 3 & 6 & 9 & 12\\
\midrule
删失人数 & 156 & 81 & 73 & 148 & 139 & 797\\
累计人数 & 156 & 237 & 310 & 458 & 597 & 1394\\
累计百分比 & 0.9 & 1.4 & 1.8 & 2.7 & 3.5 & 8.2\\
\addlinespace
住院人数 & 626 & 640 & 646 & 1265 & 1163 & 887\\
母乳喂养人数 & 15,392 & 13,128 & 10,765 & 6893 & 4717 & --\\
\bottomrule
\end{tabular}
\end{table}
\FloatBarrier

\subsection{因果假设}

为了进一步定义反事实，我们需要无个体间干扰和序贯正值性假设。
\begin{itemize}
\item \textbf{无个体间干扰（no interference）：}一名研究对象的潜在结局不依赖其他研究对象的暴露。在这里，这对应于以下假设：给定第二名婴儿的母乳喂养状态，第一名婴儿的母乳喂养不会影响第二名婴儿住院的概率。
\item \textbf{序贯正值性（sequential positivity）：}对于满足 $A_{t-1}=1$ 的每一种可能历史 $\bar L_{t-1}$，所有研究对象在时点 $t$ 继续或停止母乳喂养的概率都必须大于零。由于我们不考虑停止后又重新开始母乳喂养的策略，当 $A_{t-1}=0$ 时，不要求 $A_t$ 满足正值性。此外，对于每一种可能的历史 $\bar A_{t-1},\bar L_{t-1}$，时点 $t$ 的删失概率必须小于1。
\end{itemize}

在 PROBIT 中，无个体间干扰假设要求一名母亲的母乳喂养状态不会影响另一名母亲孩子的结局。我们认为这很合理，因为母亲在医院停留的时间较短，彼此互动有限。对于正值性，我们估计在此前持续母乳喂养的条件下继续或停止母乳喂养的概率时，没有发现接近零的取值，详见第15.3.1节。此外，在系数的所有取值下，删失概率均很低，这表明正值性在本研究中不是一个值得担忧的问题。

令 $\bar a=(a_1,a_2,\ldots,a_{K-1})$ 表示固定的母乳喂养方案。例如，母乳喂养持续超过第一个时间区间，但在第二个区间之前停止，可以写为 $(1,0,0,\ldots,0)$。由于母乳喂养近似具有单调性，我们只考虑单调的纵向暴露，即比较在不同时间停止母乳喂养的相对效应。另以 $\bar a_t=(a_1,\ldots,a_t)$ 表示这一固定方案截至时点 $t$ 的部分。

为定义目标因果参数，我们考虑一项医院层面的干预，对每名研究对象规定特定的母乳喂养持续时间。于是，对 $t>0$，可以将反事实变量 $L_t^{\bar a}$ 定义为：如果某个个体遵循所分配的母乳喂养方案 $\bar a_{t-1}$，并且保持未删失，其本会具有的观测值 $L_t$。类似地，$Y^{\bar a}$ 是在母乳喂养方案 $\bar a=\bar a_5$ 下本会观察到的反事实住院数。与这一干预对应的个体反事实数据为
\[
 O^a=(L_0,L_1^{\bar a},L_2^{\bar a},\ldots,L_{K-1}^{\bar a},Y^{\bar a}).
\]
推断目标是边际平均反事实结局，记为 $\psi_0^{\bar a}=\EE(Y^{\bar a})$。等价地，我们估计的是：如果所有婴儿均按各种不同的时间停止母乳喂养，发生住院的时间区间平均会有多少个。

为了使这一因果参数能够被估计，还需要序贯一致性和序贯可交换性（Robins，2000）。
\begin{itemize}
\item \textbf{序贯一致性（sequential consistency）：}纵向情形的一致性假设是，当 $\bar A_{t-1}=\bar a_{t-1}$ 时，$\bar L_t^{\bar a}=\bar L_t$。等价地，我们观察到的反事实变量序列，是在实际观察到的处理方案下所定义的序列。
\item \textbf{序贯可交换性（sequential exchangeability）：}给定过去的信息，反事实中间变量与最近的干预节点，即暴露和删失，相互独立：
\[
 L_t^{\bar a}\perp A_{t-1},C_t
 \mid \bar L_{t-1},\bar A_{t-2},C_{t-1}=0,
 \qquad t=1,\ldots,K+1,
\]
其中将 $\bar A_{-1}$ 视为空变量并移除。
\end{itemize}

序贯一致性假设，我们观察到的潜在结局，就是在分配某一母乳喂养持续时间这一干预下本会观察到的结局。它假设，在某个时间区间内具体何时停止母乳喂养，不影响反事实。另一种看法是，分配时允许研究对象自行决定该区间内的确切停止时间；但这种看法需要额外的可交换性条件（VanderWeele 和 Hernán，2013）。

对于上述序贯可交换性假设，必须假定 $\bar L_t$ 已经调整母乳喂养和删失的所有基线混杂因素及时变混杂因素。具体地，我们假定 $\bar L_t$ 足以控制母乳喂养 $A_t$ 的混杂，而且给定 $\bar L_{t-1}$ 后，删失 $C_t$ 可忽略。图15.1中，一次访视的删失不允许依赖最近发生的住院。虽然这一点不现实，但删失比例较低，提示这种违反可能不会严重影响分析。总体上，可交换性假设无法验证，也通常难以完全相信；不过，我们在 Schnitzer 等（2014）中认为，控制簇的指示变量可以增强这一假设的可信度。

\subsection{模型与参数}

为定义观测数据和反事实数据的模型，我们假定在患者聚集于各医院中心的总体中，所观察的医院是从某个大总体中随机抽取的。以 $O_{mi}$ 表示医院簇 $m$ 中患者 $i$ 的观测向量。第 $m$ 家医院的观测数据定义为 $O_m^c=(O_{mi};i\in Z_m)$，其中 $Z_m$ 是属于医院 $m$ 的研究对象集合。我们假定各簇观测 $O_m^c$ 从概率分布 $P_0^c$ 独立同分布抽取，并以 $O^c\sim P_0^c$ 表示这一随机变量。$O^c$ 的概率分布 $P_0^c$ 属于某个模型空间 $\mathcal M^c$。先随机选择一个医院簇，再从中随机选择一名患者，这名患者的边际概率分布 $P_0$ 对应一个混合概率分布。此边际分布可写成簇分布的函数，即 $P_0=P_0(P_0^c)$。我们直接对包含真实 $P_0$ 的边际模型空间 $\mathcal M$ 施加全部模型限制，再将各簇概率分布的模型空间限制为 $\mathcal M^c=\{P_0^c:P_0\in\mathcal M\}$，使其满足边际模型空间 $\mathcal M$ 上的约束。

在边际层面，假定 $O$ 的真实分布函数可以按个体的时间顺序分解为
\[
\begin{aligned}
 P_0={}&\underbrace{\prod_{t=1}^{K}
 Q_{0,L_t}(L_t\mid\bar C_t,\bar A_{t-1},\bar L_{t-1})Q_{0,L_0}(L_0)}_{Q_0}\\
 &\quad\times\underbrace{\prod_{t=1}^{K-1}
 G_{0,A_t}(A_t\mid\bar L_t,\bar C_t,\bar A_{t-1})
 \prod_{t=1}^{K}G_{0,C_t}(C_t\mid\bar A_{t-1},\bar L_{t-1},\bar C_{t-1})}_{G_0}.
\end{aligned}
\]
这里，$Q_0$ 是各 $L_t$ 变量的联合条件分布；$Q_{0,L_t}$（$t=1,\ldots,K$）是给定 $L_t$ 之前的信息时各 $L_t$ 的分布，$Q_{0,L_0}$ 是基线协变量的分布。类似地，$G_0$ 是暴露变量和删失变量的条件分布，可以进一步分解为各时点的分布，分别记为 $G_{0,A_t}$（$t=1,\ldots,K-1$）和 $G_{0,C_t}$（$t=1,\ldots,K$）。除对处理和删失分布 $G_0$ 施加的限制之外，模型 $\mathcal M$ 是非参数的。

在上述因果假设下，对分布 $P_0$ 应用通常的 g 公式，就可以识别目标参数 $\psi_0^{\bar a}$。设所有研究对象被分配的暴露方案固定为 $\bar a$，于是各 $A_t$ 不再是随机的；并固定所有研究对象均被完整观察，于是对 $t=1,\ldots,K$，均有 $C_t=0$。与这一静态干预 $\bar a$ 对应的边际反事实分布函数记为 $Q_0^{\bar a}$。其 g 公式为
\begin{equation}
 Q_0^{\bar a}(\bar L_K^{\bar a})=
 \prod_{t=1}^{K}Q_{0,L_t}(L_t\mid\bar C_t=0,
 \bar A_{t-1}=\bar a_{t-1},\bar L_{t-1})Q_{0,L_0}(L_0).
 \tag{15.1}
\end{equation}
我们关心的目标参数，即固定母乳喂养方案 $\bar a$ 下的边际均值，可以写为
\[
 \psi_0^{\bar a}=\Psi(P_0(P_0^c))=\Psi(P_0)
 =\Psi(Q_0^{\bar a})=\EE_{Q_0^{\bar a}}(Y^{\bar a}),
\]
其中的期望在真实的反事实数据生成函数 $Q_0^{\bar a}$ 下取得。

第15.2和15.3节估计目标参数时，暂且将不同医院的个体层面数据视为彼此独立同分布，其概率分布为模型 $\mathcal M$ 中的 $P_0$，并将目标参数作为 $P_0$ 的函数。随后，第15.4节将只承认各簇独立同分布，建立这种独立同分布型 LTMLE 的渐近线性，并提供正式推断。我们猜想，在干扰参数得到一致估计的条件下，这种独立同分布型 LTMLE 对模型 $\mathcal M^c$ 实际上也渐近有效，但本章没有对此作出正式证明。
